\bookchapter{맺음말: 호기심을 유지하고, 깨어 있으세요}

여기까지 읽으신 것을 축하합니다.

이 책은 스마트폰 하드웨어에서 소프트웨어 생태계까지, 일상 설정에서 개발자 모드까지, 보안에서 디지털 유산까지 다룹니다. 이 책이 실제 문제를 해결하는 데 도움이 되었다면 — 열지 말아야 할 권한을 끄는 법을 배우거나, 백업 습관을 들이는 것만으로도 — 이 글들이 헛되지 않은 것입니다.

하지만 어떤 구체적인 기술보다 중요한 두 가지가 있습니다:

\begin{itemize}
  \item \textbf{호기심을 유지하세요}: 스마트폰은 매년 변하고, AI도 변하고, 폴더블도 변하고, 위성 통신도 변합니다. 오늘 배운 조작 절차는 내일이면 시대에 뒤처질 수 있지만, "왜 그런지 알고 싶다"는 호기심은 절대 시대를 지나지 않습니다. "이 권한은 무엇을 하는 거지", "이 데이터는 어디에 저장되지", "이 기능이 프라이버시를 유출하지는 않을까"라는 질문을 던지는 것이 조작 절차 백 가지를 기억하는 것보다 유용합니다.
  \item \textbf{냉정함을 유지하세요}: 스마트폰은 도구이지 주인이 아닙니다. 그것은 당신의 관심을 끌고, 시간을 차지하고, 소비를 유도하도록 설계되었습니다. 사용하는 것은 당신의 선택이지만, 사용당하는 것은 당신의 통제력 상실입니다. 정기적으로 화면 시간을 올려다보고, "나는 스마트폰으로 일을 하고 있는 건지, 스마트폰에 시간을 깎이고 있는 건지" 정기적으로 생각하세요. 이것은 감상적인 것이 아니라 디지털 시대의 생존 본능입니다.
\end{itemize}

기술은 변하지만 사람의 필요는 변하지 않습니다 — 우리는 안전감, 연결, 존중이 필요합니다. 스마트폰은 그저 매개체일 뿐입니다. 진정한 “스마트폰 수업”은 책 안에 있지 않고, 매일 화면을 잠금 해제하는 그 순간에 있습니다.

\begin{tcbpractice}{마지막 연습}
이 책을 덮고, 스마트폰을 들고, 한 가지 일을 하세요:\textbf{30초 동안 생각해 보세요. 스마트폰에서 “잃어버리면 정말 아까울” 세 가지는 무엇인가요?}그리고 그것들이 백업되어 있는지 확인하세요. 이 책은 그만큼의 가치가 있습니다.

\end{tcbpractice}
